\exercisesection{2}

\begin{enumerate}
\def\labelenumi{\arabic{enumi}.}
\item
  设 K 为交换域，A 为二元不定元上的多项式环 K{[}X, Y{]}，且 \(a_{n}\) 为 A 中的主理想 \((XY^{n})\)；序列 \((a_{n})_{n\geqslant1}\) 与 \(a_{0}=A\) 一起构成 A 上的穷竭且分离滤过。设 b 为 A 的主理想 (X)；证明 A 上的拓扑在 b 上诱导出的拓扑严格粗于由 A 上的滤过在 b 上相应的滤过所定义的拓扑（从而后一个滤过当然不同于由 A 上的滤过诱导的滤过）。
\item
  设 K 为特征 \(\neq2\) 的交换域，A 为二元不定元的形式幂级数环 K{[}{[}X, Y{]}{]}。
\end{enumerate}

\begin{enumerate}
\def\labelenumi{(\alph{enumi})}
\item
  证明在 A 中主理想 \(\mathfrak{p} = (\mathrm{X}^2 -\mathrm{Y}^3)\) 是素理想。（若两个形式幂级数的乘积 \(f(\mathbf{X},\mathbf{Y})g(\mathbf{X},\mathbf{Y})\) 被 \(\mathbf{X}^2 -\mathbf{Y}^3\) 整除，先注意到 \(f(\mathrm{T}^3,\mathrm{T}^2) = 0\) 或 \(g(\mathrm{T}^3,\mathrm{T}^2) = 0\) 在形式幂级数环 \(\mathbf{K}[[\mathrm{T}]]\) 中成立；例如假设 \(f(\mathrm{T}^3,\mathrm{T}^2) = 0\)，先证明 \(f(\mathbf{X},\mathbf{Y}^2)\) 被 \(\mathbf{X} - \mathbf{Y}^3\) 以及 \(\mathbf{X} + \mathbf{Y}^3\) 整除，再考察 \(f(-\mathbf{Y}^3 +\mathbf{X},\mathbf{Y}^2)\)，最后证明 \(f(\mathbf{X},\mathbf{Y}^2)\) 被 \(\mathbf{X}^2 -\mathbf{Y}^6\) 整除。)
\item
  设 \(m\) 为 \(AX + AY\) 这一 \(A\) 的极大理想。证明 \(p\) 关于 \(m\)-进拓扑在 \(A\) 上是闭的（关于此拓扑 \(A\) 是 Hausdorff 且完备的），但在环 \(\text{gr}(A)\) 中，\(\text{gr}(p)\) 不是素理想（\(p\) 赋予由 \(A\) 上滤过诱导的滤过）。
\end{enumerate}

\begin{enumerate}
\def\labelenumi{\arabic{enumi}.}
\setcounter{enumi}{2}
\item
  设 A 为滤过环，E 为有限生成 A-模。证明若 E 赋予由 A 上滤过诱导的滤过，则 gr(E) 是有限生成的 gr(A)-模（参见习题 5 (c)）。~
\item
  给出一个双射 A-线性映射 \(u \colon \mathbf{E} \to \mathbf{F}\) 的例子，其中 \(\mathbf{E}\) 和 \(\mathbf{F}\) 是两个滤过 A-模，\(u\) 与滤过相容，但 \(\operatorname{gr}(u)\) 既非单射也非满射。（取 \(\mathbf{E} = \mathbf{A}\)、\(\mathbf{F} = \mathbf{A}\) 但赋予另一滤过，并取 \(u\) 为恒等映射。）
\item
  设 A 为带滤过 \((\mathfrak{a}_n)_{n\geqslant 0}\) 的交换滤过环，且 \(\mathfrak{a}_0 = \mathbf{A}\)。
\end{enumerate}

\begin{enumerate}
\def\labelenumi{(\alph{enumi})}
\tightlist
\item
  要使环 gr(A) 由次数有界的一族元素生成，充要条件是存在整数 q，使得对所有 \(n\)，\(\mathfrak{a}_n = \mathfrak{a}_{n+1} + \mathfrak{b}_n\)，其中 \(\mathfrak{b}_n\) 是理想 \(\mathfrak{a}_1^{\alpha_1}\mathfrak{a}_2^{\alpha_2}\ldots\mathfrak{a}_q^{\alpha_q}\) 的各项之和：
\end{enumerate}

\[
\alpha_ {i} \geqslant 0 (1 \leqslant i \leqslant q)
\]

\[
\sum_ {i = 1} ^ {q} i \alpha_ {i} = n.
\]

\[
\mathfrak {a} _ {n} = \mathfrak {a} _ {n + k} + \mathfrak {b} _ {n}
\]

\[
k > 0
\]

\begin{enumerate}
\def\labelenumi{(\alph{enumi})}
\setcounter{enumi}{1}
\item
  要使 \(\mathrm{gr}(\mathbf{A})\) 为 Noether 环，充要条件是 \(\mathbf{A} / \mathfrak{a}_1\) 为 Noether 环，(a) 的条件成立，并且对 \(i \leqslant q\)，\(\mathfrak{a}_i / \mathfrak{a}_{i+1}\) 是有限生成的 \((\mathbf{A} / \mathfrak{a}_1)\)-模（使用第 10 节定理 2 的推论 4）。
\item
  设 K 为交换域，A 为二元不定元的多项式环 K{[}X, Y{]}。在单项式 \(X^{m}Y^{n}\) 的集合上定义全序：规定当 \(X^{m}Y^{n} \leqslant X^{p}Y^{q}\)，或当 \(m + n \leqslant p + q\) 且 \(m + n = p + q\) 时，\(m \leqslant p\)。令 \((\mathbf{M}_{k})\) 为按此升序排列的 X, Y 中单项式序列，并令 \(a_{n}\) 为由指标 \(M_{k}\) 的 \(k \geqslant n\) 生成的 A 的理想。证明 \((a_{n})\) 是与 A 的环结构相容的滤过；关于此滤过，gr(A) 有零因子且不是 Noether 环，尽管 A 是 Noether 整环（使用 (b) 中的判据）；特别地，gr( \(a_{1}\) )（在 \(a_{1}\) 上取由 A 上滤过诱导的滤过）不是有限生成的 gr(A)-模，尽管 \(a_{1}\) 是有限生成的 A-模（参见 § 1, 第 2 节，命题 1 的推论）。~
\end{enumerate}

¶ 6. 设 A 为交换环，E 和 F 为两个 A-模，而 (E \(_{n}\) )（分别为 (F \(_{n}\) )）是 E（分别为 F）上由子 A-模组成的穷竭滤过。在张量积 G = E ⊗A F 上，考虑由 G \(_{n}\) = ∑ \(_{i+j=n}\) Im(E \(_{i}\) ⊗A F \(_{j}\) ) 组成的穷竭滤过。

\begin{enumerate}
\def\labelenumi{(\alph{enumi})}
\tightlist
\item
  证明复合典范同态
\end{enumerate}

\[
\begin{array}{r l}&(\mathrm{E} _ {i} / \mathrm{E} _ {i + 1}) \otimes_ {\mathrm{A}} (\mathrm{F} _ {j} / \mathrm{F} _ {j + 1}) \rightarrow\\&\qquad (\mathrm{E} _ {i} \otimes_ {\mathrm{A}} \mathrm{F} _ {j}) / (\operatorname{Im} (\mathrm{E} _ {i} \otimes_ {\mathrm{A}} \mathrm{F} _ {j + 1}) + \operatorname{Im} (\mathrm{E} _ {i + 1} \otimes_ {\mathrm{A}} \mathrm{F} _ {j})) \rightarrow \mathrm{G} _ {i + j} / \mathrm{G} _ {i + j + 1}\end{array}
\]

是一个次数为 0 的分次同态（称为典范同态）\(\operatorname{gr}(\mathbf{E}) \otimes_{\mathbf{A}} \operatorname{gr}(\mathbf{F}) \to \operatorname{gr}(\mathbf{E} \otimes_{\mathbf{A}} \mathbf{F})\) 的限制，该同态为满射。

\begin{enumerate}
\def\labelenumi{(\alph{enumi})}
\setcounter{enumi}{1}
\item
  证明若 \(\operatorname{gr}(\mathbf{E})\) 是平坦 A-模，则当 \(\mathrm{E}_m / \mathrm{E}_n\) 时 A-模 \(m \leqslant n\) 以及当 \(\mathrm{E} / \mathrm{E}_n\) 时 \(n \in \mathbf{Z}\) 都是平坦的。
\item
  以下假设 \(\mathrm{gr}(\mathbf{E})\) 是平坦 A-模，且 \(\mathbf{E}\) 是平坦 A-模（若 \((\mathbf{E}_n)\) 是离散滤过，第二个假设由第一个假设推出）。于是 \(\mathrm{H}_{ij} = \mathrm{E}_i \otimes_{\mathbf{A}} \mathrm{F}_j\) 被典范地识别为 \(\mathbf{G} = \mathbf{E} \otimes_{\mathbf{A}} \mathbf{F}\) 的子模（第一章，第 2 节，第 5 节，命题 4）。证明对于 \(\mathbf{R}\) 的任意两个有限子集 \(\mathbf{S}\)、\(\mathbf{Z} \times \mathbf{Z}\)，
\end{enumerate}

\[
\left(\sum_ {(i, j) \in \mathbf {R}} \mathrm{H} _ {i j}\right) \cap \left(\sum_ {(h, k) \in \mathbf {S}} \mathrm{H} _ {h k}\right) = \mathrm{H} _ {\sup (i, h), \sup (j, k)}\tag{*}
\]

其中第二个和式中的 \((i,j,h,k)\) 遍历 \(R \times S\)。（若 R 和 S 各自只含一个元素，使用第一章，第 2 节，第 6 节，命题 7。若 p 是 R 或 S 的元素中出现的指标 i 和 h 的最大值，q 是其中出现的指标 j 和 k 的最小值，考察 (*) 两边在典范同态

\[
\mathbf {E} \otimes_ {\mathbf {A}} \mathbf {F} _ {q} \rightarrow (\mathrm{E} / \mathrm{E} _ {p}) \otimes_ {\mathbf {A}} \mathbf {F} _ {q}
\]

下的像，并对 \(\operatorname{Card}(\mathbf{R}) + \operatorname{Card}(\mathbf{S})\) 作归纳。）

\begin{enumerate}
\def\labelenumi{(\alph{enumi})}
\setcounter{enumi}{3}
\tightlist
\item
  由 (c) 推出 (a) 中定义的典范同态于是为双射。
\end{enumerate}

\begin{enumerate}
\def\labelenumi{\arabic{enumi}.}
\setcounter{enumi}{6}
\tightlist
\item
  设 A 为交换环，E 为 A-模，F 为 E 的子模，B 为外代数 \(\wedge\) (E)，而 \(\Im\) 为 B 中由 F 的元素在 B 中的典范像生成的双边理想。令 \(\operatorname{gr}^{\Im}(B)\) 为 B 按 \(\Im\)-进滤过所得的伴随分次环。
\end{enumerate}

\begin{enumerate}
\def\labelenumi{(\alph{enumi})}
\item
   定义一个满射的典范（非分次）A-代数同态 \((\wedge (\mathbf{F})) \otimes_{\mathbf{A}}^{g}(\wedge (\mathbf{E} / \mathbf{F})) \to \mathrm{gr}^{\Im} (\mathbf{B})\)（这里指分次代数的斜张量积（代数，第 III 章）。
\item
  证明若 \(\mathbf{F}\) 在 \(\mathbf{E}\) 中有补子模，则 (a) 中定义的同态是同构。
\item
   取 \(\mathbf{A} = \mathbf{Z}\)、\(\mathbf{E} = \mathbf{Z} / 4\mathbf{Z}\)、\(\mathbf{F} = 2\mathbf{Z} / 4\mathbf{Z}\)；证明此时 (a) 中定义的同态不是单射。
\end{enumerate}

\begin{enumerate}
\def\labelenumi{\arabic{enumi}.}
\setcounter{enumi}{7}
\item
  设 A 为滤过环，\((\mathbf{A}_{n})\) 为其滤过，E 为滤过 A-模，\((\mathbf{E}_{n})\) 为其滤过。证明若 \(A_{0} = A\) 且 \(E_{0} = E\)，则从 \((a, x) \mapsto ax\) 到 E 的映射 \(A \times E\) 关于由这些滤过定义的拓扑是一致连续的。
\item
  给出两个滤过环 A、B 的例子，它们的滤过是穷竭且分离的，并给出一个与滤过相容的非满射同态 \(u: A \to B\)，使得 \(\text{gr}(u)\) 为双射。由此在不假设环 A 完备时，给出命题 12 的推论 1 的反例；并在不再假设 E 有限生成时，给出命题 13 的反例（使用代数，第 VIII 章，第 7 节，习题 3(b)）。
\item
  设 A 为 Noether 环，\(\sigma\) 为 A 的自同构。证明代数，第 IV 章，第 5 节，习题 10(b) 中定义的环 E 是（左或右）Noether 环。
\item
  证明若 E 是交换域上维数 \(\geqslant2\) 的向量空间，则 E 的张量代数是既非左 Noether 环也非右 Noether 环的环（若 a、b 是 E 中两个线性无关的向量，考察由 \(a^{n}b^{n}\) 时的元素 \(n\geqslant1\) 生成的左理想（或右理想））。
\end{enumerate}

 ¶ 12. 设 K 为交换域，A 为任意无限族不定元上的多项式环 \(K[X_{i}]_{i\in I}\)，m 为 A 中由 \(X_{i}\) 生成的（极大）理想。若令 \(A_{i}=A/m^{i+1}\)，则 \(A_{i}\) 以及当 \(h_{ij}:A/m^{j+1}\to A/m^{i+1}\) 时的典范同态 \(i\leqslant j\) 满足命题

14；环 \(\lim A_{i}\) 是完备化 \(\hat{A}\)（关于 m-进拓扑的 \(A\)），且典范同态 \(\hat{A} \to A_{i}\) 的核等于 \((m^{i})^{\wedge}\)，即 \(m^{i}\) 在 \(\hat{A}\) 中的闭包。

\begin{enumerate}
\def\labelenumi{(\alph{enumi})}
\item
  证明 \(\hat{\mathbf{A}}\) 被典范地识别为关于 \(\mathbf{X}_{\iota}\) 的形式幂级数环，其中每个幂级数在给定次数下只有有限项（代数，第 IV 章，第 5 节，习题 1）。
\item
  从现在起取 \(\mathbf{I} = \mathbf{N}\)。证明 \(\hat{\mathfrak{m}} \neq \hat{\mathbf{A}}. \mathfrak{m}\)（考察形式幂级数 \(\sum_{n=1}^{\infty} \mathbf{X}_n^n\)）。
\item
  设 K 为有限域。证明 \((\hat{\mathfrak{m}})^{2} \neq (\mathfrak{m}^{2})^{\bullet}\)。（先证明如下结果：对每个整数 k \textgreater{} 0，存在整数 \(n_{k}\)，并且对每个整数 \(n \geqslant n_{k}\)，存在一个以 K 为系数、在 \(F_{n}\) 个不定元中次数为 n 的齐次多项式 \(n^{2}\)，使得 \(F_{n}\) 不可能是形如 \(P_{1}Q_{1} + \cdots + P_{k}Q_{k}\) 的任一多项式的 n 次项之和，其中 \(P_{i}\) 和 \(Q_{i}\) 是同一组 \(n^{2}\) 个不定元中的无常数项多项式。）
\end{enumerate}

推出 \(\hat{A}\) 关于 \(\hat{m}\)-进拓扑不是完备的。

\begin{enumerate}
\def\labelenumi{\arabic{enumi}.}
\setcounter{enumi}{12}
\tightlist
\item
  设 K 为域，\(A = K[[X]]\) 为形式幂级数环，m 为其极大理想，故 A 关于 m-进拓扑是 Hausdorff 且完备的（第 6 节，命题 6 的推论）。在加法群 A 上考虑滤过 \((\mathrm{E}_{n})\)，其中 \(E_{0} = A\)，而 \(E_{n}\) 是 \(m^{n}\) 与环 K{[}X{]} 的交；此滤过是穷竭且分离的，它在 A 上定义的拓扑 T 与加法群结构相容并且比 m-进拓扑更细，但 A 关于拓扑 T 不是完备群（考察多项式序列 \((1 - X^{n})/(1 - X)\)）。
\end{enumerate}

14 设 A 为环（不必交换），E 为左 A-模。若 E 上的拓扑在平移下不变且 0 有一个由 E 的子模组成的邻域基本系，则称该拓扑为线性拓扑；此时称 E 带线性拓扑。线性拓扑与 E 的加法群结构相容，并且当 A 取离散拓扑时在 E 上定义拓扑 A-模结构。任意 A-模上的离散拓扑和最粗拓扑都是线性拓扑。

\begin{enumerate}
\def\labelenumi{(\alph{enumi})}
\item
  若 E 是带线性拓扑的模，F 是 E 的子模，则 F 上的诱导拓扑以及 E 关于 F 的商拓扑都是线性拓扑。若 \((\mathrm{E}_{\alpha}, f_{\alpha\beta})\) 是带线性拓扑的 A-模的逆系统，其中 \(f_{\alpha\beta}\) 是连续线性映射，则拓扑 A-模 \(E = \lim_{\leftarrow} E_{\alpha}\) 带线性拓扑。
\item
  设 E 为带线性拓扑的 A-模。E 中存在一个由开（也是闭的）子模组成的 0 的邻域基本系 \((\mathrm{V}_{\lambda})_{\lambda \in L}\)；若 \(\phi_{\lambda\mu}: E/V_{\mu} \to E/V_{\lambda}\) 为典范映射且 \(V_{\lambda} \supset V_{\mu}\)，则族 \((\mathrm{E}/\mathrm{V}_{\lambda}, \phi_{\lambda\mu})\) 是离散 A-模的逆系统；拓扑 A-模 \(\tilde{E} = \lim_{\leftarrow} E/V_{\lambda}\) 被识别为 E 的 Hausdorff 完备化，因而其带线性拓扑（一般拓扑，第 III 章，第 7 节，第 3 节）。
\item
  Hausdorff 的带线性拓扑 A-模 E 若为 Artin 模，或 E 的子模 \(\neq0\) 的集合中存在最小元，则称其为离散的。
\end{enumerate}

\begin{enumerate}
\def\labelenumi{\arabic{enumi}.}
\setcounter{enumi}{14}
\tightlist
\item
  \begin{enumerate}
  \def\labelenumii{(\alph{enumii})}
  \tightlist
  \item
    设 E 为带线性拓扑的 A-模（习题 14）。由 E 中线性簇组成的滤基具有聚点的充要条件是存在一个包含它的、由仿射线性簇组成的收敛滤基 \(B' \supset B\)。若 E 是 Hausdorff 的，且 E 上每个由仿射线性簇组成的滤基都至少有一个聚点，则称 E 是线性紧的。每个 Artin 模在离散拓扑下都是线性紧的。Hausdorff 的带线性拓扑模 F 的每个线性紧子模在 F 中都是闭的。
  \end{enumerate}
\end{enumerate}

\begin{enumerate}
\def\labelenumi{(\alph{enumi})}
\setcounter{enumi}{1}
\item
  若 E 是线性紧 A-模，且 u 是从 E 到 Hausdorff 的带线性拓扑 A-模 F 的连续线性映射，则 \(u(\mathbf{F})\) 是 F 的线性紧子模。
\item
  设 E 为 Hausdorff 的带线性拓扑 A-模，F 为 E 的闭子模。E 线性紧的充要条件是 F 和 E/F 都线性紧。
\item
  线性紧模的任意乘积都是线性紧的（考察在线性仿射簇滤基中极大的滤基。
\item
  设 \((\mathrm{E}_{\alpha}, f_{\alpha \beta})\) 是关于有向指标集的带线性拓扑模的逆系统；假设 \(f_{\alpha \beta}\) 是连续线性映射，且当 \(\alpha \leqslant \beta\) 时，\(f_{\alpha \beta}^{-1}(0)\) 是 \(\mathbf{E}_{\beta}\) 的线性紧子模。令 \(\mathbf{E} = \varprojlim \mathbf{E}_{\alpha}\)，令 \(f_{\alpha}\) 为典范映射 \(\mathbf{E} \to \mathbf{E}_{\alpha}\)；证明对所有 \(\alpha, f_{\alpha}(\mathbf{E}) = \bigcap_{\alpha \leqslant \beta} f_{\alpha \beta}(\mathbf{E}_{\beta})\)（特别地，若 \(f_{\alpha \beta}\) 为满射，则 \(f_{\alpha}\) 也是满射），且 \(f_{\alpha}^{-1}(0)\) 是线性紧的（使用 (d) 以及一般拓扑，第 I 章，附录，第 2 节，定理 1）。
\end{enumerate}

¶ 16. (a) 设 E 为 Hausdorff 的带线性拓扑 A-模。证明下列性质等价：

(α) E 是线性紧的（习题 15）。

(β) 对每个连续线性映射 \(u\)，它从 \(E\) 映到 Hausdorff 的带线性拓扑 A-模 \(F\)，\(u(E)\) 都是 \(F\) 的闭子模。

\begin{enumerate}
\def\labelenumi{(\Alph{enumi})}
\setcounter{enumi}{24}
\tightlist
\item[(γ)]
  对 E 上每个比给定拓扑粗的线性拓扑（无论是否 Hausdorff），E 都是完备的。
\end{enumerate}

(δ) E 是完备的，且 E 中存在一个由子模组成的 0 的开邻域基本系 \((\mathbf{U}_{\lambda})\)，使得 \(E/U_{\lambda}\) 是离散的线性紧 A-模（参见 § 3，习题 5）。~

（为看出 \((\gamma)\) 蕴含 \((\delta)\)，取 \(F\) 为 \(E\) 的一个开子模，以及由仿射线性簇组成的滤基 \(\mathfrak{B}\) 在 \(E / F\) 上。对所有 \(V \in \mathfrak{B}\)，令 \(M_V\) 为 V 在 E/F 中的方向子模在 E 中的逆像；在 E 上考虑以 \(M_{V}\) 为 0 的邻域基本系的线性拓扑。）

\begin{enumerate}
\def\labelenumi{(\alph{enumi})}
\setcounter{enumi}{1}
\item
  设 E 为 Hausdorff 的带线性拓扑 A-模，M 为 E 的线性紧子模。证明对于 E 的每个闭子模 F，\(M + F\) 在 E 中是闭的（考察 M 在 E/F 中的像）。
\item
  设 E 为线性紧 A-模，F 为 Hausdorff 的带线性拓扑 A-模。证明 \(E \times F\) 的每个闭子模 M 在 F 上的投影都是闭的。得到逆命题（参见一般拓扑，第 I 章，第 10 节，第 2 节）。~
\item
  设 E 为线性紧 A-模，u 为从 E 到 Hausdorff 的带线性拓扑 A-模 F 的连续线性映射。证明对于 E 上每个由仿射线性簇组成的滤基 B，B 的聚点集在 u 下的像就是 \(u(\mathfrak{B})\) 的聚点集。特别地，对于 E 的每个闭子模 M，都有 \(\bigcap_{N\in\mathfrak{B}}(M+N)=M+\bigcap_{N\in\mathfrak{B}}N\)。
\end{enumerate}

\begin{enumerate}
\def\labelenumi{\arabic{enumi}.}
\setcounter{enumi}{16}
\tightlist
\item
  若 A-模 E 上的线性拓扑 \(\mathcal{T}\) 是 Hausdorff 的，且不存在严格粗于 \(\mathcal{T}\) 的 Hausdorff 线性拓扑，则称其为极小的。
\end{enumerate}

\begin{enumerate}
\def\labelenumi{(\alph{enumi})}
\item
  Hausdorff 线性拓扑 T 在 E 上为极小的充要条件是：E 上每个由仿射线性簇组成且只有一个聚点的滤基 B 都收敛到该点。（为看出条件是必要的，注意当 M 遍历 B、V 遍历由子模组成的 0 的邻域基本系时，M + V 构成一个与 B 有相同聚点的滤基。为看出条件是充分的，注意由开子模组成且交约化为 0 的滤基，是比 T 粗的某个 Hausdorff 线性拓扑中 0 的邻域基本系。）
\item
  A-模 E 上的离散拓扑为极小的充要条件是 E 为 Artin 模。
\item
  若 \(\mathcal{T}\) 为极小拓扑，则 \(\mathcal{T}\) 在 \(\mathbf{E}\) 的任意闭子模上诱导的拓扑也是极小的。
\end{enumerate}

¶ 18. 在 A-模 E 中，若子模 M ≠ E 且 E/M 的非零子模集合中存在最小元，则称 M 为受保护的。

\begin{enumerate}
\def\labelenumi{(\alph{enumi})}
\item
  证明 E 的每个子模 \(N \neq E\) 都是若干受保护子模的交（对于所有 \(x \notin N\)，考察 E 中包含 N 但不包含 x 的子模集合中的极大元）。
\item
  设 \(\mathcal{T}\) 是 E 上的 Hausdorff 线性拓扑，令 \(\mathcal{T}^*\) 为如下线性拓扑：其 0 的邻域基本系是由 E 中关于 \(\mathcal{T}\) 开且受保护的子模生成的滤基。证明 \(\mathcal{T}^*\) 是 Hausdorff 的，且关于 \(\mathcal{T}\) 闭的每个子模关于 \(\mathcal{T}^*\) 也闭（注意关于 \(\mathcal{T}\) 闭的每个子模都是关于 \(\mathcal{T}\) 开的子模的交）。推出若 E 关于 \(\mathcal{T}^*\) 完备，则 E 关于 \(\mathcal{T}\) 完备（一般拓扑，第 III 章，第 3 节，第 5 节，命题 9）。
\item
  假设 \(\mathcal{T}\) 是线性紧的。证明 \(\mathcal{T}^*\) 是线性紧的，并且是比 \(\mathcal{T}\) 粗的 Hausdorff 线性拓扑中最粗的；特别地，\(\mathcal{T}^*\) 是极小线性拓扑（习题 17）。（令 \(\mathfrak{B}\) 为由关于 \(\mathcal{T}\) 闭的子模组成且交为 0 的滤基；使用习题 16 (d)，证明对于关于 \(\mathcal{T}\) 开且受保护的每个子模 U，存在 \(\mathbf{M} \in \mathfrak{B}\) 使得 \(\mathbf{M} \subset \mathbf{U}\)。）
\item
  设 \(\mathbf{F}\) 为 Hausdorff 的带线性拓扑 A-模，\(\mathcal{T}_1\) 为其拓扑；证明若 \(u: \mathrm{E} \to \mathrm{F}\) 关于拓扑 \(\mathcal{T}\) 和 \(\mathcal{T}_1\) 是连续线性映射，则 \(u\) 关于拓扑 \(\mathcal{T}^*\) 和 \(\mathcal{T}_1^*\) 也连续。
\end{enumerate}

¶ 19. (a) 设 E 为线性紧 A-模。证明下列条件等价：

(α) 对从 E 到 Hausdorff 的带线性拓扑 A-模 F 的每个连续线性映射 u，u 都是从 E 到 F 的严格态射（一般拓扑，第 III 章，第 2 节，第 8 节）。

(β) 对 \(\mathbf{F}\) 的每个闭子模 \(\mathbf{E}\)，\(\mathbf{E} / \mathbf{F}\) 上的商拓扑都是极小拓扑（习题 17）。

(γ) E 是完备的，且 E 中存在一个由子模组成的 0 的开邻域基本系 \((\mathbf{U}_{\lambda})\)，使得 \(E/U_{\lambda}\) 是 Artin 模。

（为看出 \((\gamma)\) 蕴含 \((\beta)\)，将其化约到 \(F = 0\) 的情形，并使用习题 17 和 18。）

若 E 满足等价条件 \((\alpha)\)、\((\beta)\) 和 \((\gamma)\)，则称 E 为严格线性紧的。

\begin{enumerate}
\def\labelenumi{(\alph{enumi})}
\setcounter{enumi}{1}
\item
  设 E 为 Hausdorff 的带线性拓扑 A-模，F 为 E 的闭子模。E 严格线性紧的充要条件是 F 和 E/F 都严格线性紧。
\item
  严格线性紧模的每个逆极限都是严格线性紧的。
\item
  设 E 为 Hausdorff 的带线性拓扑 A-模，u 为从 E 到严格线性紧 A-模 F 的线性映射。证明若 u 的图像在 \(E \times F\) 中闭，则 u 连续（使用习题 16 (c) 以及如下事实：若 M 在 E 中闭且 E/M 是 Artin 模，则 M 在 E 中开）。
\end{enumerate}

¶ 20. (a) 证明离散的线性紧模不可能是无穷多个不约化为 0 的子模的直和。

\begin{enumerate}
\def\labelenumi{(\alph{enumi})}
\setcounter{enumi}{1}
\item
  给出一个不是线性紧的极小线性拓扑（习题 17）的例子（考察无穷个两两不同构的简单模的直和）。
\item
  设 E 为线性紧模，且存在一个简单子模族，其和在 E 中稠密。证明：(1) E 严格线性紧（将 (a) 应用于 E 对一个开子模的商）；(2) E 同构于一族离散简单模的（拓扑）乘积。（考察 E 的极大开子模集合 O，使得对于每个有限
\end{enumerate}

\[
\mathrm{E} / \left(\bigcap_ {k = 1} ^ {n} \mathrm{G} _ {k}\right)
\]

 个由序列 \((\mathbf{G}_k)_{1\leqslant k\leqslant n}\) 中的 \(n\) 个属于 \(\mathfrak{O}\) 的互异元素组成的序列，\(\operatorname{E} / \left(\bigcap_{k=1}^{\infty}\mathbf{G}_k\right)\) 是 \(n\) 个简单子模的直和。证明存在一个（关于包含关系）极大集合 \(\mathfrak{O}\)，且属于此集合的所有子模 \(\mathbf{G}\) 的交约化为 0；使用如下事实：\(\{0\}\) 是极大开子模的交，并且对于每个极大开子模 \(\mathbf{G}\)，它是 \(\mathbf{E}\) 的子模，\(\mathbf{E}\) 至少有一个不包含于 \(\mathbf{G}\) 的简单子模。最后使用 \(\mathbf{E}\) 严格线性紧这一事实得出结论。）

\begin{enumerate}
\def\labelenumi{(\alph{enumi})}
\setcounter{enumi}{3}
\tightlist
\item
  由 (c) 推出每个域 \(\mathbf{K}\) 上的线性紧向量空间都是严格线性紧的，并同构于乘积 \(\mathbf{K}_s^{\mathrm{I}}\)。
\end{enumerate}

¶ 21. 回忆一下，在环 A 上，若一个拓扑是 A-模 \(\mathbf{A}_s\) 上的线性拓扑且与 A 的环结构相容，则称其为（左）线性的（第二章，第 2 节，习题 16）。若拓扑环 A 的 \(\mathbf{A}_s\) 是线性紧（分别严格线性紧）A-模，则称 A 为（左）线性紧（分别（左）严格线性紧）的。

\begin{enumerate}
\def\labelenumi{(\alph{enumi})}
\item
  证明若 A 关于拓扑 T 是左线性紧的，则习题 18 (b) 中定义的拓扑 \(T^{*}\) 也与 A 的环结构相容（使用习题 18 (d)）。
\item
  假设 \(\mathbf{A}\) 是交换的；令 \(u \neq 0\) 为 \(\mathbf{A} / \Re\) 的幂等元，其中 \(\Re\) 是 \(\mathbf{A}\) 的 Jacobson 根。证明若 \(\mathbf{A}\) 取离散拓扑且为线性紧的，则存在唯一的幂等元 \(e \in \mathbf{A}\)，其在 \(\mathbf{A} / \Re\) 中的像等于 \(u\)。（在线性仿射簇 \(x + \mathfrak{b}\) 中，其中 \(\mathfrak{b} \subset \Re\) 且 \(x^2 - x \in \mathfrak{b}\)，且该簇包含于类 \(u\)，证明存在最小元 \(e + \mathfrak{a}\)。使用代数，第 VIII 章，第 6 节，习题 10(a)，推出 \(e^2 = e\) 且 \(\mathfrak{a} = 0\)。考察 \(e^2 - e = r\)，证明 \(r \in \mathrm{Ar}^2\)，并证明 \(e\) 的唯一性：若 \(e_1, e_2\) 是 \(\mathbf{A}\) 的两个幂等元且 \(e_1 e_2 \in \Re\)，则 \(e_1 e_2 = 0\)。）
\item
  证明每个取离散拓扑且线性紧的交换环 A 都是有限个局部环的直积（这些局部环也取离散拓扑且线性紧）。（先考虑 \(\Re = 0\) 的情形，使用习题 15 (b) 以及第二章，第 1 节，第 2 节，命题 5；然后将 (b) 应用于 A/ \(\Re\) 的幂等元。）
\item
  证明每个线性紧（分别严格线性紧）交换环 A 都是线性紧（分别严格线性紧）局部环族的乘积。（令 \((\mathfrak{m}_{\lambda})\) 为 A 的开极大理想族；对于每个开理想 \(a\)，它不包含 \(\mathfrak{m}_{\lambda}\)，令 \(\mathrm{A}_{\lambda}^{\prime} \times \mathrm{A}_{\lambda}^{\prime \prime}\) 为 A/a 作为两个环的直积分解，使 \(\mathrm{A}_{\lambda}^{\prime}\) 是 (c) 中定义的、极大理想为 \((\mathfrak{m}_{\lambda} + \mathfrak{a}) / a\) 的局部环；令 \(e_{\lambda}^{\prime}(a)\) 为 \(\mathrm{A}_{\lambda}^{\prime}\) 的单位元，并将其视为 A 中模 \(a\) 的类；\(e_{\lambda}^{\prime}(a)\) 构成一个在 A 中收敛到幂等元 \(e_{\lambda}\) 的滤基，且 \(Ae_{\lambda}\) 是 A 的闭理想，\(e_{\lambda}\) 是其单位元。证明环 \(Ae_{\lambda}\) 是局部环，且 A 同构于各个 \(Ae_{\lambda}\) 的乘积。）
\end{enumerate}

¶ 22. (a) 设 A 为局部环，m 为其极大理想。证明若关于 A 上的线性拓扑 T，A 是严格线性紧的，则 T 比 m-进拓扑粗。

\begin{enumerate}
\def\labelenumi{(\alph{enumi})}
\setcounter{enumi}{1}
\item
  设 A 为交换环，m 为 A 的理想。A 关于 m-进拓扑严格线性紧的充要条件是：A 关于此拓扑 Hausdorff 且完备，A/m 是 Artin 环，且 \(m/m^{2}\) 是一个有限长度的 (A/m)-模 *（换言之，A 是完备的 Noether 半局部环）*。
\item
  设 I 为无限指标集，K 为有限域，\(A = K[[X_{t}]]_{t \in I}\) 为关于不定元族 \((\mathbf{X}_{t})_{t \in I}\) 的形式幂级数代数（代数，第 IV 章，第 5 节，习题 1）；A 是局部环。作为 K 上的向量空间，A 可识别为乘积空间 \(\mathbf{K}^{\mathbf{N}^{(1)}}\)；若 K 取离散拓扑而 A 取乘积拓扑 T，证明 T 是环 A 上的线性拓扑，且 A 关于此拓扑严格线性紧；若 m 是 A 的极大理想，仿照习题 12 (c) 证明 A 关于 m-进拓扑不是完备的。
\end{enumerate}

¶ 23. 设 A 为带线性拓扑 T 的交换环。

\begin{enumerate}
\def\labelenumi{(\alph{enumi})}
\item
  证明 \(\mathfrak{a}\) 的如下理想 \(\mathbf{A}\)：它们关于 \(\mathcal{T}\) 是闭的且 \(\mathbf{A} / \mathfrak{a}\) 是 Artin 环，构成 \(\mathcal{T}^{(c)}\) 上某个线性拓扑 \(\mathbf{A}\) 的 0 的邻域基本系，而该拓扑比 \(\mathcal{T}\) 粗（使用代数，第 VIII 章，第 2 节，习题 12）；然后 \((\mathcal{T}^{(c)})^{(c)} = \mathcal{T}^{(c)}\)。令 \(\mathbf{B}\) 为 \(\mathbf{A}\) 关于拓扑 \(\mathcal{T}^{(c)}\) 的 Hausdorff 完备环；证明 \(\mathbf{B}\) 是严格线性紧的；称 \(\mathbf{B}\) 为与 \(\mathbf{A}\) 相伴的严格线性紧环。
\item
  令 \(\mathcal{T}_{\omega}(\mathbf{A})\) 为 \(\mathbf{A}\) 上的离散拓扑，\(\mathcal{T}_c(\mathbf{A})\) 为拓扑 \((\mathcal{T}_{\omega}(\mathbf{A}))^{(c)}\)。拓扑 \(\mathcal{T}^{(c)}\) 在 \(\mathbf{A}\) 上为 Hausdorff 的充要条件是 \(\mathcal{T}\) 为 Hausdorff，且对于每个开理想 \(\mathfrak{b}\)，它是 \(\mathbf{A}\) 中关于 \(\mathcal{T}\) 开的理想，拓扑 \(\mathcal{T}_c(\mathbf{A}/\mathfrak{b})\) 为 Hausdorff。* 若 \(\mathbf{A}\) 是高度 1 的非离散赋值的环（第 VI 章），则 \(\mathcal{T}_c(\mathbf{A})\) 不是 Hausdorff。*
\item
  假设 A 关于 \(\mathcal{T}\) 是线性紧的；则 A 关于 \(\mathcal{T}^{(c)}\) 是完备的（但一般不是 Hausdorff 的）。\(\mathcal{T}^{(c)}\) 为 Hausdorff 的充要条件是 A 上存在一个严格线性紧且比 \(\mathcal{T}\) 粗的线性拓扑；此时 \(\mathcal{T}^{(c)}\) 是具有这些性质的唯一拓扑。
\item
  设 E 为带线性拓扑且严格线性紧的 A-模。证明若 A 取拓扑 \(\mathcal{T}_{c}(A)\)，则 E 是拓扑环（注意若 F 是 Artin A-模，则对所有 \(x \in F\)，x 的湮灭子 a 使得 A/a 是 Artin 环）。推出若 B 是 A 关于拓扑 \(\mathcal{T}_{c}(A)\) 的 Hausdorff 完备化，则 E 可看作拓扑 B-模；环 B 同构于严格线性紧局部环的乘积
\end{enumerate}

\(B_{\lambda}\)（习题 21 (d)），证明 E 同构于严格线性紧子模 \(E_{\lambda}\) 的乘积，其中 \(E_{\lambda}\) 被指标 \(B_{\mu}\) 的 \(\mu \neq \lambda\) 湮灭，因而可看作拓扑 \(B_{\lambda}\)-模（若 \(e_{\lambda}\) 是 B 的单位元，取 \(E = e_{\lambda}.E\)）。

\begin{enumerate}
\def\labelenumi{\arabic{enumi}.}
\setcounter{enumi}{23}
\tightlist
\item
  在交换环 A 上，令 \(\mathcal{T}_{m}(A)\)~ 表示如下线性拓扑：其 0 的邻域基本系由有限个极大理想的幂的乘积（等于交）组成（参见第 13 节，命题 17）；若 A 中任意有限个理想 \(\neq0\) 的交都是 \(\neq0\)，令 \(\mathcal{T}_{u}(A)\) 表示如下线性拓扑：其 0 的邻域基本系由 A 的所有理想 \(\neq\{0\}\) 组成。
\end{enumerate}

\begin{enumerate}
\def\labelenumi{(\alph{enumi})}
\item
  证明 \(\mathcal{T}_c(\mathbf{A})\)（习题 23）比 \(\mathcal{T}_m(\mathbf{A})\) 粗（注意 Artin 环的 Jacobson 根是幂零的）；给出 \(\mathcal{T}_c(\mathbf{A}) \neq \mathcal{T}_m(\mathbf{A})\) 的例子（参见习题 22 (c)）。~* 若 A 是高度 1 的非离散赋值的环（第 VI 章），则 \(\mathcal{T}_c(\mathbf{A}) = \mathcal{T}_m(\mathbf{A})\)，而 \(\mathcal{T}_u(\mathbf{A}) \neq \mathcal{T}_m(\mathbf{A})\)。*
\item
  拓扑 \(\mathcal{T}_m(A)\) 是 \(A\) 上满足以下条件的最粗线性拓扑：\(A\) 的极大理想都是闭的，且每个开理想的每个幂都是开理想（注意在 \(A\) 上使极大理想 \(m\) 闭的线性拓扑下，\(m\) 必然是开的）。
\item
  若 A 是 Noether 环，则 \(\mathcal{T}_{m}(\mathbf{A}) = \mathcal{T}_{c}(\mathbf{A})\)（注意若 \(m\) 是 A 的极大理想，则对每个整数 \(\mathbf{A} / m^{k}\)，\(k \geqslant 1\) 是 Artin 环，并注意 \(\mathfrak{m}^h / \mathfrak{m}^{h+1}\) 是必然具有有限长度的 (A/m)-模）。给出一个满足 \(\mathcal{T}_{m}(\mathbf{A}) \neq \mathcal{T}_{u}(\mathbf{A})\) 的 Noether 局部环 A 的例子（考察域上多项式环的一个局部环）。
\end{enumerate}

\begin{enumerate}
\def\labelenumi{\arabic{enumi}.}
\setcounter{enumi}{24}
\item
  设 A 为拓扑环，E 为有限生成左 A-模。证明 E 上存在一个拓扑（称为典范拓扑），该拓扑与其 A-模结构相容且比其他所有拓扑都细（将 E 写成 \(A_{s}^{n}/R\) 的形式，在 \(A_{s}^{n}\) 上取乘积拓扑，在 E 上取商拓扑）。若 \(E'\) 是拓扑 A-模且 \(u: E \to E'\) 是 A-线性映射，证明 u 关于 E 上的典范拓扑是连续的。若进一步 \(E'\) 有限生成且带典范拓扑，则 u 是严格态射。若 A 上的拓扑由滤过（\(A_{n}\)）定义，证明 E 上的典范拓扑由滤过（\(A_{n}E\)）定义。
\item
  设 A 为交换环，m 为 A 的理想，\((a_{\lambda})_{\lambda \in L}\) 为 m 的一个生成系，\(\hat{A}\) 为 A 关于 m-进拓扑的 Hausdorff 完备化。
\end{enumerate}

 证明若 \(\mathbf{A}'\) 表示关于不定元族 \((\mathrm{T}_{\lambda})_{\lambda \in \mathbf{L}}\) 的形式幂级数环，且每个给定次数只有有限项（习题 12），则 \(\hat{\mathbf{A}}\) 同构于 \(\mathbf{A}'\) 对理想的闭包 \(\mathfrak{b}\) 的商，其中该理想是 \(\mathbf{A}'\) 中由 \(T_{\lambda} - a_{\lambda}\) 生成的理想（使用第 8 节定理 1）。特别地，p-进整数环 \(Z_{p}\) 同构于商环 \(\mathbf{Z}[[T]]/(T-p)\)。

¶ 27. (a) 设 A 为带线性拓扑的交换环。对于 A 的每个乘性子集 S，令 \(\mathrm{A}\{\mathrm{S}^{-1}\}\) 表示环 \(\mathrm{S}^{-1}\mathrm{A}\) 的 Hausdorff 完备化，其拓扑的 0 的邻域基本系由理想 \(\mathrm{S}^{-1}\mathrm{U}_{\lambda}\) 组成，其中 \((\mathrm{U}_{\lambda})\) 是 A 中由 A 的理想组成的 0 的邻域基本系。证明 \(\mathrm{A}\{\mathrm{S}^{-1}\}\) 典范同构于环 \(\mathrm{S}_{\lambda}^{-1}(\mathrm{A}/\mathrm{U}_{\lambda})\) 的逆极限，其中 \(\mathrm{S}_{\lambda}\) 是 S 在 \(\mathrm{A}/\mathrm{U}_{\lambda}\) 中的典范像。若 \(\hat{\mathrm{A}}\) 是 A 的 Hausdorff 完备化，则 \(\mathrm{A}\{\mathrm{S}^{-1}\}\) 典范同构于 \(\hat{\mathrm{A}}\{\mathrm{S}'^{-1}\}\)，其中 \(\mathrm{S}'\) 是 S 在 \(\hat{\mathrm{A}}\) 中的典范像。

\begin{enumerate}
\def\labelenumi{(\alph{enumi})}
\setcounter{enumi}{1}
\item
  \(A\{S^{-1}\}\) 约化为 0 的充要条件是：在 A 中，0 属于 S 的闭包。举出一个 A 是 Hausdorff 且完备、但 \(S^{-1}A\) 关于 (a) 中定义的拓扑并非如此的例子。
\item
  对于每个连续同态 \(u\)，它从 \(A\) 到完备 Hausdorff 的带线性拓扑环 \(B\)，若 \(u(S)\) 由 \(B\) 的可逆元组成，证明 \(u = u' \circ j\)，其中 \(j: A \to A\{S^{-1}\}\) 是典范映射且 \(u'\) 连续；此外 \(u'\) 唯一确定。
\item
  设 \(S_{1}\)、\(S_{2}\) 是 A 的两个乘性子集，\(S_{2}^{\prime}\) 是 \(S_{2}\) 在 \(A\{S^{-1}\}\) 中的典范像；定义 \(A\{(S_{1}S_{2})^{-1}\}\) 到 \(A\{S_{1}^{-1}\}\{S_{2}'^{-1}\}\) 的典范同构。
\item
  设 \(\mathfrak{a}\) 为 A 的开理想，\(\mathfrak{a}\{\mathrm{S}^{-1}\}\) 为 \(\mathrm{S}^{-1}\mathfrak{a}\) 关于 \(\mathrm{S}^{-1}\mathrm{A}\) 上诱导拓扑的 Hausdorff 完备化；证明 \(\mathfrak{a}\{\mathrm{S}^{-1}\}\) 被典范地识别为 \(\mathrm{A}\{\mathrm{S}^{-1}\}\) 的一个开理想，且离散环 \(\mathrm{A}\{\mathrm{S}^{-1}\} / \mathfrak{a}\{\mathrm{S}^{-1}\}\) 同构于 \(\mathrm{S}^{-1}(\mathrm{A} / \mathfrak{a})\)。反之，若 \(\mathfrak{a}'\) 是 \(\mathrm{A}\{\mathrm{S}^{-1}\}\) 的开理想，则其在 A 中的逆像 \(\mathfrak{a}\) 是一个开理想，且 \(\mathfrak{a}' = \mathfrak{a}\{\mathrm{S}^{-1}\}\)。特别地，映射 \(\mathfrak{p} \mapsto \mathfrak{p}\{\mathrm{S}^{-1}\}\) 将 A 中不与 S 相交的开素理想集递增双射到 \(\mathrm{A}\{\mathrm{S}^{-1}\}\) 的开素理想集。
\item
  设 \(p\) 为 \(A\) 的开素理想，令 \(S = A - p\)。证明 \(A\{S^{-1}\}\) 是局部环，其剩余域同构于 \(A / p\) 的分式域。
\item
  对于每个元素 \(f \in \mathbf{A}\)，令 \(\mathbf{A}_{(f)}\) 表示环 \(\mathbf{A}\{\mathbf{S}_f^{-1}\}\)，其中 \(\mathbf{S}_f\) 是由 \(f^n\)（\(n \geqslant 0\)）组成的乘性集。若 \(f\) 遍历乘性子集 \(\mathbf{S}\)（属于 \(\mathbf{A}\)），则 \(\mathbf{A}_{(f)}\) 构成一个环的直接系统，其直接极限记为 \(\mathbf{A}_{(\mathbf{S})}\)（不带任何拓扑）；定义一个典范同态
\end{enumerate}

\[
\mathrm{A} _ {\{\mathrm{S} \}} \rightarrow \mathrm{A} \{\mathrm{S} ^ {- 1} \}.
\]

 若 \(S = A - p\)，其中 \(p\) 是 \(A\) 的开素理想，证明 \(A_{(S)}\) 是局部环，典范同态 \(A_{(S)} \to A\{S^{-1}\}\) 是局部同态，且 \(A_{(S)}\) 与 \(A\{S^{-1}\}\) 的剩余域典范同构（参见第二章，第 3 节，习题 16）。~

\begin{enumerate}
\def\labelenumi{(\alph{enumi})}
\setcounter{enumi}{7}
\tightlist
\item
   假设 A 上的拓扑是某个 A 的理想 m 的 m-进拓扑，A 是 Hausdorff 且完备的，且 \(m/m^{2}\) 是有限生成的 (A/m)-模。证明若 \(m' = m\{S^{-1}\}\)，则 \(A' = A\{S^{-1}\}\) 上的拓扑是 \(m'\)-进拓扑，\(m' = mA'\)，且 \(m'/m'^{2}\) 是有限生成的 \(A'/m'\)-模（使用第 10 节命题 14）。若 A 是 Noether 环，则 \(A'\) 也是 Noether 环。
\end{enumerate}

\begin{enumerate}
\def\labelenumi{\arabic{enumi}.}
\setcounter{enumi}{27}
\tightlist
\item
  设 A 为带线性拓扑的交换环，E、F 为带线性拓扑的两个拓扑 A-模。若 V（分别 W）遍历 E（分别 F）的开子模集合，则子模
\end{enumerate}

\[
\operatorname{Im} (\mathrm{V} \otimes_ {\mathbf {A}} \mathrm{F}) + \operatorname{Im} (\mathrm{E} \otimes_ {\mathbf {A}} \mathrm{W})
\]

在 \(\mathbf{E} \otimes_{\mathbf{A}} \mathbf{F}\) 中构成某个 0 的邻域基本系，该 \(\mathbf{E} \otimes_{\mathbf{A}} \mathbf{F}\) 的拓扑与其在拓扑环 A 上的模结构相容，称为 E 和 F 上给定拓扑的张量积。此 A-模的 Hausdorff 完备化 \((\mathbf{E} \otimes_{\mathbf{A}} \mathbf{F})^{\wedge}\) 是一个 \(\hat{\mathbf{A}}\)-模，称为 E 与 F 的完备张量积。

\begin{enumerate}
\def\labelenumi{(\alph{enumi})}
\tightlist
\item
  证明若 \((\mathrm{V}_{\lambda})\)（分别 \((\mathrm{W}_{\mu})\)）是 E（分别 F）中由子模组成的 0 的邻域基本系，则 \((\mathbf{E} \otimes_{\mathbf{A}} \mathbf{F})^{\wedge}\) 典范同构于如下 A-模逆系统的逆极限
\end{enumerate}

\[
\left(\mathrm{E} / \mathrm{V} _ {\lambda}\right) \otimes_ {\mathrm{A}} \left(\mathrm{F} / \mathrm{W} _ {\mu}\right);
\]

推出 \((\mathbf{E} \otimes_{\mathbf{A}} \mathbf{F})^{\wedge}\) 是一个 \(\hat{\mathbf{A}}\)-模，且典范同构于 \((\hat{\mathbf{E}} \otimes_{\hat{\mathbf{A}}} \hat{\mathbf{F}})\)；它也记作 \(\hat{\mathbf{E}} \widehat{\otimes}_{\hat{\mathbf{A}}} \hat{\mathbf{F}}\)。

\begin{enumerate}
\def\labelenumi{(\alph{enumi})}
\setcounter{enumi}{1}
\item
  设 \(\mathbf{E}'\)、\(\mathbf{F}'\) 为两个带线性拓扑的拓扑 A-模，且 \(u: \mathbf{E} \to \mathbf{E}', v: \mathbf{F} \to \mathbf{F}'\) 为两个连续 A-线性映射；证明 \(u \otimes v: \mathbf{E} \otimes \mathbf{F} \to \mathbf{E}' \otimes \mathbf{F}'\) 关于 \(\mathbf{E} \otimes \mathbf{F}\) 和 \(\mathbf{E}' \otimes \mathbf{F}'\) 上的张量积拓扑是连续的；将连续线性映射记为 \(u \hat{\otimes} v\)，其为 \((\mathbf{E} \otimes \mathbf{F})^{\wedge} \to (\mathbf{E}' \otimes \mathbf{F}')^{\wedge}\) 且对应于 \(u \otimes v\)。
\item
  设 B、C 为带线性拓扑的两个交换 A-代数，使得典范映射 \(\mathbf{A} \to \mathbf{B}\)、\(\mathbf{A} \to \mathbf{C}\) 连续（为简便起见，称 B、C 为拓扑 A-代数）。证明 \((\mathbf{B} \otimes_{\mathbf{A}} \mathbf{C})^{\wedge}\) 具有典范的拓扑 A-代数结构，称为代数 B、C 的完备张量积。定义典范连续表示
\end{enumerate}

\[
\rho \colon \mathbf {B} \rightarrow (\mathbf {B} \otimes_ {\mathbf {A}} \mathbf {C}) ^ {\wedge}, \quad \sigma \colon \mathbf {C} \rightarrow (\mathbf {B} \otimes_ {\mathbf {A}} \mathbf {C}) ^ {\wedge}
\]

具有如下性质：对于每个完备 Hausdorff 交换拓扑 A-代数 D 以及每一对连续 A-同态 \(u\colon B\to D\)、\(v\colon C\to D\)，存在唯一的连续 A-同态 \(w\colon(B\otimes_{A}C)^{\wedge}\to D\)，使得 \(u=w\circ\rho\) 且 \(v=w\circ\sigma\)。

\begin{enumerate}
\def\labelenumi{(\alph{enumi})}
\setcounter{enumi}{3}
\tightlist
\item
  设 m 为 A 的理想，且 A 上的拓扑是 m-进拓扑。若 E 和 F 都赋予 m-进拓扑，证明在 \(E \otimes_{A} F\) 上，E 和 F 上拓扑的张量积就是 m-进拓扑。推出 \((\mathbf{E} \otimes_{\mathbf{A}} \mathbf{F})^{\wedge}\) 同构于逆极限 \(\varprojlim ((\mathrm{E} / \mathfrak{m}^{n+1}) \otimes_{\mathbf{A}} \mathbf{F})\) 和 \(\varprojlim (\mathbf{E} \otimes_{\mathbf{A}} (\mathbf{F} / \mathfrak{m}^{n+1}\mathbf{F}))\)。
\end{enumerate}

\begin{enumerate}
\def\labelenumi{\arabic{enumi}.}
\setcounter{enumi}{28}
\item
  设 A 为交换环，m 为 A 的理想，且 \(\bigcap_{n\geq0}m^{n}=0\)，c 为 A 中不是零因子的元素。证明传递子 \(q_{n}=m^{n}:Ac\) 关于 m-进拓扑是开的，且其交为 0。若 A 关于 m-进拓扑严格线性紧（习题 22(b)），则 \((q_{n})\) 是此拓扑下 A 中 0 的邻域基本系。
\item
  设 K 为无限交换域，A 为二元不定元的形式幂级数环 \(K[[X,Y]]\)，它是完备的 Hausdorff Noether 局部环。定义 A 的一个乘性子集 S，使得 \(S^{-1}A\) 不是半局部环（若 \((\lambda_{n})\) 是 K 中互异元素的无限序列，考察 A 的素理想 \(p_{n}=A(X+\lambda_{n}Y)\)）。
\item
  证明若 A 是半局部环，则形式幂级数环 A{[}{[}X{]}{]} 也是半局部环（考察此环的 Jacobson 根）。
\end{enumerate}

